\documentclass[11pt,twoside,fontset=none,scheme=plain]{ctexart}

\usepackage[
  paperwidth=176mm,
  paperheight=250mm,
  top=21mm,
  bottom=22mm,
  inner=23mm,
  outer=19mm,
  headheight=14pt,
  headsep=7mm,
  footskip=13mm
]{geometry}
\usepackage[nomath,semibold]{libertinus}
\usepackage{microtype}
\usepackage{xcolor}
\usepackage{hyperref}
\usepackage{fancyhdr}

\newcommand{\writingfontpath}{assets/fonts/}
\IfFileExists{assets/fonts/LXGWWenKaiGB-Regular.ttf}{}{%
  \renewcommand{\writingfontpath}{../../assets/fonts/}%
}
\setCJKmainfont[
  Path=\writingfontpath,
  BoldFont=LXGWWenKaiGB-Medium.ttf,
  AutoFakeSlant=0.2
]{LXGWWenKaiGB-Regular.ttf}
\setCJKsansfont[
  Path=\writingfontpath,
  BoldFont=LXGWWenKaiGB-Medium.ttf
]{LXGWWenKaiGB-Medium.ttf}
\setCJKmonofont[
  Path=\writingfontpath,
  BoldFont=LXGWWenKaiGB-Medium.ttf
]{LXGWWenKaiGB-Regular.ttf}
\xeCJKsetup{PunctStyle=kaiming}

\definecolor{bookink}{HTML}{222222}
\definecolor{mutedink}{HTML}{6B6964}
\definecolor{accent}{HTML}{315F86}
\newcommand{\writingtwodigits}[1]{\ifnum#1<10 0\fi\number#1}
\newcommand{\writinglastrevised}{\number\year-\writingtwodigits{\month}-\writingtwodigits{\day}}
\hypersetup{
  unicode=true,
  pdftitle={关于我：世界，爱情与意识形态},
  pdfauthor={Xayah Hina},
  pdfdisplaydoctitle=true,
  colorlinks=true,
  linkcolor=accent,
  urlcolor=accent,
  citecolor=accent
}
\urlstyle{same}

\setlength{\parindent}{2em}
\setlength{\parskip}{0pt}
\linespread{1.28}
\emergencystretch=1.5em
\clubpenalty=10000
\widowpenalty=10000
\displaywidowpenalty=10000
\setcounter{secnumdepth}{3}
\raggedbottom

\ctexset{
  section={
    name={},
    number=\arabic{section},
    format=\Large\sffamily\bfseries\color{bookink},
    aftername=\hspace{0.8em},
    beforeskip=2.6ex plus 0.5ex minus 0.2ex,
    afterskip=1.3ex plus 0.2ex,
    indent=0pt
  },
  subsection={
    name={},
    number=\thesection.\arabic{subsection},
    format=\large\sffamily\bfseries\color{bookink},
    aftername=\hspace{0.75em},
    beforeskip=2.2ex plus 0.4ex minus 0.2ex,
    afterskip=1ex plus 0.2ex,
    indent=0pt
  },
  subsubsection={
    name={},
    number=\thesubsection.\arabic{subsubsection},
    format=\normalsize\sffamily\bfseries\color{bookink},
    aftername=\hspace{0.7em},
    beforeskip=1.8ex plus 0.3ex minus 0.2ex,
    afterskip=0.8ex plus 0.2ex,
    indent=0pt
  }
}

\renewcommand{\abstractname}{摘要}
\renewenvironment{abstract}{%
  \begin{center}
    \small\sffamily\bfseries\color{bookink}\abstractname
  \end{center}
  \begin{list}{}{\leftmargin=2em\rightmargin=2em}
    \item\relax\small\color{mutedink}\noindent\ignorespaces
}{%
  \end{list}
  \vspace{1.2em}
}

\makeatletter
\renewcommand{\maketitle}{%
  \thispagestyle{plain}%
  \begin{center}
    \vspace*{0.6em}
    {\LARGE\sffamily\bfseries\color{bookink}\@title\par}
    \vspace{1em}
    {\small\color{mutedink}\@author\quad\textperiodcentered\quad\@date\par}
  \end{center}
  \vspace{1.8em}
}
\makeatother

\pagestyle{fancy}
\fancyhf{}
\fancyhead[LE]{\footnotesize\color{mutedink}Xayah Hina}
\fancyhead[RO]{\footnotesize\color{mutedink}关于我：世界，爱情与意识形态}
\fancyfoot[LE,RO]{\footnotesize\color{mutedink}\thepage}
\renewcommand{\headrulewidth}{0pt}
\fancypagestyle{plain}{%
  \fancyhf{}
  \fancyfoot[C]{\footnotesize\color{mutedink}\thepage}
  \renewcommand{\headrulewidth}{0pt}
}
\AtBeginDocument{\color{bookink}}

\title{关于我：世界，爱情与意识形态}
\author{Xayah Hina}
\date{2026-07-17\quad\textperiodcentered\quad Last Revised: \writinglastrevised}

\begin{document}
\maketitle

\begin{abstract}
诚实地直面自己，记录此刻幼稚的思维。以自己的历史为鉴，为今后优化自己提供solid baseline。
\end{abstract}
% Begin writing here.

\section{写在前面}
本文中的内容可以预见，在一年后甚至三个月后的自己看来都极其浅薄，贻笑大方。因此权且当作一个naive的baseline，用于记录一下自己浅薄无知的伤春悲秋时代的认知。


\section{世界、历史、政治与“我”}
自从我开始正式理解整个世界开始，无时无刻不处在世界与“我”的巨大矛盾之中。我该怎么认识这个世界？我在这个世界存在的意义是什么？这不是教科书中的哲学思辨，而是贯穿整个建构出“我”的意识形态所要回答的重要问题。

认识世界，最重要的一点莫过于认识国家认识政府。正确认识自己的国家与政府，这个对你有生杀予夺管辖权的庞然大物让人不得不直面它正视它。正确认识别的国家，是一种朴素认知下的乖离，一种允许超脱规则禁忌下直面真实世界的途径。

\subsection{如何看待中国与中国共产党}
\begin{quote}
\textit{这个小标题我写下来都觉得有些心惊。但我完全不认同墙内社区自我审查的习惯，我认为语言就要保持纯净表达，力求直接、简明扼要，不要拐弯抹角。}
\end{quote}

这是出生、生活在这片土地的人不得不一生直面的问题。究竟应该怎么看待中国和中国共产党？是正义的还是邪恶的？应该合作还是逃离？哪些地方是好的哪些地方是坏的？中国是什么？中国共产党又是什么？

接下来我决定在正式向自己回答一下在此刻（2026年）的答案。

中国，对于我而言自始至终都是一个美好的概念。960万平方公里的土地尤其是故乡的土地每一寸都令人眷恋，这一点在每一次留学回国落地之后最为强烈。中国对于我来说是一种文化根植，是一种逃避不了的深刻思想烙印。所以我说这份感情是“自始至终”的。对于中国的历史，也是精彩与悲恸交织，令人感慨万千。溢美之词其实在本文意义不大所以暂且不多言，于我而言，对于中国、文化、语言这类我自始至终都是高度评价的。

下一个就是比较敏感的问题：如何看待中国共产党？这个问题也是必须回答而且必须深思熟虑后回答的。这里对过去的认知不多花篇幅来讲，重点说当前的认知。中国共产党在历史上无疑是需要被归类于邪恶与极端一类的，这一点自正式理解历史之后我的认知是一贯的。但有没有真的那么十恶不赦呢？从客观上并没有。如果从十分评价制来评价我个人心目中，1949\textasciitilde 1957为8/10分。从某种意义上来说，共产党从事实上清除了沉疴弊病，巨大的群众动员能力破除了许多民国乃至封建时代的痼疾，并且防止了民族分裂割据，而且稳定了国际局势。从某种意义上来说，当时的领导人确实可以被称作拥有卓越的能力。1958\textasciitilde 1976为4/10分。这一点也是这一年我认知变化比较大的一点。这二十年，混乱与恐惧确实笼罩了历史上整一代人，但真的仔细去认真研究历史文献之后，倒是让我有了一定的改观。并且从历史进程来看，混乱也并非一无是处，有很多被动带来的益处，这不详细赘述。至于1977\textasciitilde 至今，如今的我可以评价为9.5/10分。经济领域取得的成就有目共睹自然不必多说，政治体制过去我多有微词，但如果真正静下心诘问自己，事实上我是完全认同当前的政治制度的，而且不是基本认同而是完全认同，无论是选举制度还是政府组织制度。

\subsection{如何看待日韩台、欧洲和美国}
时间走到2026年，随着我走过和认识的其他国家越来越多，对于其他国家的认知也在逐步变化。但毕竟没有深入考察过大部分国家的真实情况，本节能写的不是很多。

对于日本，这是一个美丽的国家，但总感觉有点死气沉沉没有生气。如果跳脱中国人的身份站在完全客观的视角来看，日本国在历史上是一个强大的国家这点毋庸置疑。我自己本身就是认可强大的力量的人，无论所谓正义还是邪恶。因此自明治维新依赖这个国家一直都处于强大之列，哪怕二战之后也处于。但近些年来日本似乎受到日本文化自身的强大桎梏，并且缺乏新鲜血液，显得充满了暮气。但需要承认，日本终是一个强大的、与中国文化最为接近的发达国家，因此我不得不认识它、直面它、并且尽可能保持足够的联系。

对于韩国我的了解不多，但是韩国在这些年在我心中的好感度陡升。需要承认，至少从文化层面，韩国有许多值得了解和学习的地方。

对于台湾，我的印象颇为不佳。继承了中华民国的衣钵，近些年来越来越让人难以理解。甚至对于在国外遇到的台湾的民众，在交流时很明显能感受到一种对大陆人敏感、抵触和防备。

对于欧洲，我的印象也是非常差。和日本一样充满暮气，但又多了一种完全异域的隔阂感。欧洲离巴尔干、北非、印度太近了，自由主义下社会治理越来越混乱，让人担忧。

最后是一个不得不面对的国家，唯一的超级大国美国。对于美国的感情我是复杂的。说起来也是愚钝，我对于美国强大的首次感受是在本科的选修课上无意看到的slide上的一张表格，上面写了美国的贫困率统计标准，这给了懵懂无知的我极大震撼。在这之前我我知道美国富裕，强大，也知道中国最顶尖的人才都会去美国深造并且以留在美国拿绿卡换护照为最终目标，但对于究竟富裕成什么样子完全没有一个概念。对于美国，我想是每一个想要向上爬的人永远无法绕过的一个国家。对于我而言，美国不是一个国家，不像日本/韩国/中国/英国这样的民族国家。美国更像是一个副本，就像游戏中最终挑战boss的那个地狱副本。美国是一个移民国家，这就意味着事实上没有一个确切的“美国人”的概念。换句话说，你也可以是美国人，我也可以是美国人。第一次意识到这个问题的时候给了我很大的震撼，按照这样思考之后让我视野忽然扩大了很多，也忽然想明白了很多问题。对于今天的我而言，美国是一个极难的副本，一旦决定进入就要面临极大的挑战，有可能遇到极大的危险葬送自己，也有可能淘到金实现美国梦。从某种意义上来说，我认为美国至今为止还是一个勇敢者的副本游戏。

\subsection{中国、世界与我}

最后，一切都必须回到“我”上面来。那么，“我”该怎么办？我要怎么处理自己的世界观和方法论？

从某种意义上来说，我是一个保守主义者。我喜欢凸性机会，但又极度恐惧风险。在本节开始我提到了，国家与政府，是对你有生杀予夺管辖权的庞然大物。因此，无论如何，保全自己，保全家人和爱人都是对我而言的第一任务。如果规避风险，如何及时脱离险境是我一生中必须时刻准备的课题。

说实话，当前中国体制下，法律自然是不可靠的。因此想要保全自己，最有效的方法是制衡。或者换句话说，让自己有统战价值。很多时候，你不需要直面最高权力和国家机器，你需要处理的只是在社区、区县一级的权力，因此你必须有能力做到这一层级的制衡。制衡不意味着要显式抛头露面，最好的制衡是不露声色的。古话有言：一等洋人二等官，三等少民四等汉。对于这个体制，即使我对它高度评价，但也不得不提防，因此这二等之路这辈子也是不打算趟一遭浑水了。遍观当前整个精英阶级，一等路线时至今日仍然是最大的一条正确路径。但对我而言，这并不意味着就必须要移民或者长居国外，而是必须让自己有底牌，时刻有fallback的最终选择。毕竟对于我而言，我需要的只是提防而不是厌恶。如果长久在一个更加差的环境下度过余生，反而是得不偿失。毕竟，我认为我始终是一个保守主义者，我最大的愿望是保全自己，保全家人和爱人，抵御风险，安稳过日子。


\section{合作，爱情与女性权利}
上一节讨论了在社会中生存的“我”对于社会组织的形式“国家”和“政府”的思辨。在本章中，我计划讨论一下对于在社会中生存的“我”对于人际关系形式的“爱情”的思辨。

为什么要谈论爱情。很不幸，这是作为生物底层写就的烙印，作为人类始终无法回避这个问题。无论如何，自己都不得不承认：人是社会性的，人是有欲望的。这里的欲望不完全指向性，或者说我更希望把它定义为爱。

我必须承认，人是需要被爱的。

TODO：我还没想好怎么写

\end{document}
